\section{Materials and Methods}
\label{sec:methods}

The proposed study was designed to evaluate whether the elementary fault
mechanisms forming a compound motor fault can be recovered when their exact
combination has not been observed during model development. Rather than
treating each compound condition as an independent multiclass label, the
problem was formulated at the constituent-mechanism level. The methodology
consisted of a primary crossed constituent-generalization experiment
together with complementary analyses designed to distinguish diagnostic
failure from absence, suppression, or context-dependent reorientation of
constituent evidence. The analysis pipeline included constituent-level
fault representation, physics-guided multimodal feature construction,
independent constituent detection, crossed evaluation of composition,
compound-context, and operating-condition generalization, within-pair and
cross-partner transferability analysis, physical-signature preservation
analysis, and paired statistical evaluation. All preprocessing, feature
normalization, classifier fitting, and decision-threshold calibration were
performed exclusively within the corresponding training data to prevent
information leakage.

\subsection{Constituent-Level Fault Representation and Problem Formulation}
\label{sec:representation}

Let the dataset contain $N$ independent experimental runs. Instead of
assigning each run to one of the original fault-condition classes, every
run was represented by a binary constituent vector
\begin{equation}
\mathbf{y}_i = \left[y_{i,1}, y_{i,2}, \ldots, y_{i,K}\right]^{T},
\qquad \mathbf{y}_i \in \{0,1\}^{K},
\label{eq:convec}
\end{equation}
where $i$ denotes the experimental run and $K = 9$ denotes the number of
elementary fault mechanisms considered in this study. The constituent set
was
\begin{equation}
\mathcal{F} = \{f_{\mathrm{ball}}, f_{\mathrm{inner}}, f_{\mathrm{outer}},
f_{\mathrm{bend}}, f_{\mathrm{bar}}, f_{\mathrm{dyn}}, f_{\mathrm{stat}},
f_{\mathrm{vu}}, f_{\mathrm{wind}}\},
\label{eq:conset}
\end{equation}
corresponding to bearing-ball, bearing-inner-race, bearing-outer-race,
shaft-bending, broken-bar, dynamic-eccentricity, static-eccentricity,
voltage-unbalance, and winding-related fault mechanisms, respectively.

For a healthy run,
\begin{equation}
\lVert \mathbf{y}_i \rVert_0 = 0,
\end{equation}
whereas a single-fault run satisfies
\begin{equation}
\lVert \mathbf{y}_i \rVert_0 = 1.
\end{equation}
A compound fault is represented by the simultaneous activation of two or
more constituents,
\begin{equation}
\lVert \mathbf{y}_i \rVert_0 \geq 2.
\end{equation}
The composition associated with the $i$th run was therefore defined as
\begin{equation}
\mathcal{C}_i = \left\{ f_k \in \mathcal{F} : y_{i,k} = 1 \right\}.
\end{equation}
For example, a simultaneous inner-bearing and winding fault was encoded as
\begin{equation}
y_{i,\mathrm{inner}} = 1, \qquad y_{i,\mathrm{wind}} = 1,
\end{equation}
while all remaining constituent entries were zero.

This formulation differs fundamentally from conventional closed-set
multiclass diagnosis. In a multiclass formulation, a previously unseen
compound such as $A{+}B$ constitutes an unseen class. In the proposed
constituent formulation, the diagnostic objective is instead to recover the
known elementary mechanisms $A$ and $B$ independently. Consequently, the
primary research question becomes whether constituent knowledge learned
from single faults and other compound contexts can be transferred to a
novel composition.

\subsection{Physics-Guided Multimodal Feature Representation}
\label{sec:features}

The diagnostic representation was derived from three vibration channels and
three phase-current channels sampled at 12{,}800~Hz. Each 2-s analysis
window was angle resampled using the measured keyphase pulses at 256
samples per shaft revolution before order-domain feature extraction.
Window-level descriptors were subsequently aggregated to one run-level
representation using the median across valid windows. Interquartile-range
descriptors were also generated during feature extraction, but the curated
mechanism-specific signatures used in the final constituent detectors
excluded them. Thus, all classifiers operated on one physics-guided feature
vector per experimental run. The resulting multimodal representation was
written as
\begin{equation}
\mathbf{x}_i = \left[\left(\mathbf{x}_i^{(v)}\right)^{T},
\left(\mathbf{x}_i^{(c)}\right)^{T}\right]^{T},
\end{equation}
where $\mathbf{x}_i^{(v)}$ and $\mathbf{x}_i^{(c)}$ denote vibration- and
current-derived descriptors, respectively.

For an angle-domain signal $x[n]$, the order spectrum $P_x(o)$ was
estimated using a Hann-window periodogram. Around a target shaft order
$o_0$, local spectral energy was accumulated within a $\pm 0.12$ order band
and normalized by the total spectral energy between 0.25 and 60 shaft
orders,
\begin{equation}
R_x(o_0) =
\frac{\displaystyle\sum_{|o-o_0| \le 0.12} P_x(o)}
     {\displaystyle\sum_{0.25 \le o \le 60} P_x(o) + \epsilon}.
\label{eq:orderratio}
\end{equation}
Both the local band energy and its logarithmic normalized ratio,
$\log_{10}\!\left[\max\left(R_x(o_0), \epsilon\right)\right]$, were
retained in the feature bank.

For the SKF~6205 bearing, the characteristic shaft orders were
\begin{equation}
o_{\mathrm{BPFO}} = 3.585, \qquad
o_{\mathrm{BPFI}} = 5.415, \qquad
o_{\mathrm{BSF}} = 2.357,
\label{eq:orders}
\end{equation}
together with their second harmonics. Bearing-fault descriptors were
extracted from the Hilbert-envelope vibration order spectrum.

Current-domain descriptors were obtained from the individual phase-current
order spectra and the magnitude of the Clarke space vector. The dominant
electrical order was identified from the maximum summed three-phase
spectral energy over 0.5--30 orders. Relative to this data-derived carrier
$o_e$, features were evaluated at
\begin{equation}
o_e, \qquad o_e \pm 1, \qquad o_e \pm 2, \qquad 2 o_e.
\label{eq:sidebands}
\end{equation}
The $o_e \pm 1$ and $o_e \pm 2$ components were used as electrical-carrier
sideband proxies and were not interpreted as a fully slip-parameterized
broken-bar model.

Current imbalance was characterized using the coefficient of variation of
the three phase RMS values and a phase-order-invariant sequence ratio. Let
$S_1$ and $S_2$ denote the two sequence candidates obtained from the three
complex phase-current coefficients at $o_e$. Because the dataset channel
phase order was not assumed \emph{a priori}, the implemented sequence
descriptor was
\begin{equation}
r_{\mathrm{seq}} =
\frac{\min\left(|S_1|, |S_2|\right)}
     {\max\left(|S_1|, |S_2|\right) + \epsilon}.
\label{eq:seqratio}
\end{equation}
Accordingly, this descriptor represents a phase-order-invariant sequence
imbalance rather than the magnitude of a prespecified negative-sequence
phasor.

Rather than exposing every constituent detector to the complete feature
bank, a fixed mechanism-specific mapping was used. Bearing-inner,
bearing-outer, and bearing-ball detectors used BPFI/2BPFI, BPFO/2BPFO, and
BSF/2BSF envelope-order descriptors, respectively. Shaft bending used raw
shaft-$1\times$ and shaft-$2\times$ vibration features. Broken-bar
detection used electrical-carrier sideband proxies; dynamic and static
eccentricity used combined shaft-order and carrier-sideband descriptors;
voltage unbalance used sequence-ratio and phase-RMS imbalance features; and
winding faults used sequence-ratio, phase-RMS imbalance, and electrical
second-harmonic descriptors. For constituent $k$,
\begin{equation}
\mathbf{x}_i^{(k)} = \mathbf{M}_k \mathbf{x}_i,
\label{eq:selector}
\end{equation}
where $\mathbf{M}_k$ denotes the predefined mechanism-specific feature
selector. The selected sets contained 16 descriptors for each bearing
mechanism and shaft bending, 40 for broken bar, 56 for each eccentricity
mechanism, 2 for voltage unbalance, and 12 for winding. The unequal
selector sizes reflect the number of physically distinct channel--order
combinations represented by each predefined signature family rather than
post-hoc optimization of feature-set dimensionality. No outer-test
performance was used to modify the feature definitions or selector sizes.

\subsection{Independent Constituent Detection}
\label{sec:detection}

A separate binary detector was trained for each constituent mechanism,
allowing multiple fault mechanisms to be activated simultaneously without a
softmax constraint. For the locked primary analysis, each detector used
$L_2$-regularized logistic regression with $C = 1$ and class balancing
applied within the corresponding training data. The predefined
mechanism-specific feature subset $\mathbf{x}_i^{(k)}$ described in
Section~\ref{sec:features} was used independently for constituent $k$.

All preprocessing was fitted strictly within each outer-training fold.
Missing values were replaced using training-set medians, after which
features were standardized using training-set statistics. The resulting
probability $p_{i,k}$ was converted to a binary decision according to
\begin{equation}
\hat{y}_{i,k} = \mathbb{I}\left(p_{i,k} \ge \tau_k\right),
\label{eq:threshold}
\end{equation}
where $\tau_k$ is a constituent-specific threshold.

Thresholds were calibrated exclusively within the outer-training data using
grouped inner cross-validation over operating conditions. Candidate
thresholds from 0.05 to 0.95 were evaluated using constituent-level $F_1$
as the primary objective, with balanced accuracy and proximity to 0.5 used
as deterministic tie breakers. The selected threshold was then fixed for
the corresponding outer-test fold. No outer-test sample was used for
imputation, normalization, classifier fitting, or threshold selection.

A deliberately simple linear classifier was retained as the locked primary
diagnostic probe because the objective was to characterize constituent
transferability rather than to optimize model architecture. This choice
limits architecture-induced confounding when comparing compound-context
regimes. Alternative classifier families were evaluated separately as
post-hoc robustness analyses.

\subsection{Severity-Merging Sensitivity Analysis}
\label{sec:severity}

The locked mechanism-level formulation merges high- and low-severity
single-fault examples belonging to the same constituent. Because the
released compound states predominantly contain the high-severity variants
of the severity-coded mechanisms, a post-hoc sensitivity analysis was
performed to determine whether low-severity single-fault support materially
affected the crossed generalization results. The analysis removed the
low-severity single-fault runs for bearing inner, bearing outer, static
eccentricity, and winding, corresponding to 48 runs in total. The 108
compound test runs were left unchanged. The same constituent
representation, curated feature sets, logistic-regression classifier,
crossed outer splits, and grouped inner-CV threshold-calibration procedure
were then repeated on the reduced 240-run training-support dataset.
Differences from the locked H/L-merged primary analysis were evaluated
using 5000 paired physical-run cluster-bootstrap resamples. This experiment
was treated only as a sensitivity analysis and was not used to redefine the
locked primary method.

\subsection{Classifier-Family Robustness Analysis}
\label{sec:classifierfamilies}

To determine whether the observed constituent-generalization patterns were
specific to the locked linear diagnostic probe, the complete crossed
protocol was repeated with three additional classifier families: a random
forest, an RBF-kernel support-vector machine, and a three-hidden-layer
feed-forward neural network (MLP). These analyses were performed as
post-hoc robustness experiments and were not used to replace or select the
locked logistic-regression primary model.

The random forest used 100 trees, square-root feature subsampling, a
minimum leaf size of two, and balanced-subsample class weighting. The
RBF-SVM used $C = 1$, gamma set to \texttt{scale}, and balanced class
weighting. The neural baseline used hidden-layer sizes of 64, 32, and 16
units with ReLU activation, Adam optimization, $L_2$ regularization of
$10^{-4}$, an initial learning rate of $10^{-3}$, and a maximum of 400
iterations. Class balancing for the neural baseline was implemented using
balanced sample weights when supported and deterministic balanced
oversampling otherwise.

Across classifier families, the constituent labels, curated
mechanism-specific feature sets, crossed outer splits, training-only median
imputation and standardization, grouped inner cross-validation, threshold
search grid, threshold-selection objective, and evaluation metrics were
held fixed. No alternative classifier or hyperparameter configuration was
selected on the basis of outer-test performance.

\subsection{Closed-Set Control Experiment}
\label{sec:closedset}

An auxiliary closed-set experiment was conducted to establish the
conventional diagnostic separability of the dataset before imposing
compositional generalization constraints. The original diagnostic states
were treated as 24 mutually exclusive classes and evaluated using five-fold
cross-validation at the experimental-run level. All windows or derived
descriptors belonging to the same physical run were assigned to the same
fold. Preprocessing parameters were estimated exclusively from the training
runs of each fold, and performance was calculated from out-of-fold
predictions so that each physical run contributed to the test set exactly
once. This experiment was used only as a closed-set reference and was not
used for model selection in the constituent-generalization experiments.

\subsection{Crossed Composition--Context--Condition Generalization
Protocol}
\label{sec:protocol}

To separate the effects of novel fault combinations, compound-context
exposure, and operating-condition shift, a crossed outer evaluation
protocol was constructed over all available compound runs. The compound
structure can equivalently be represented as an undirected graph
\begin{equation}
\mathcal{G} = (\mathcal{V}, \mathcal{E}),
\label{eq:graph}
\end{equation}
where the vertex set $\mathcal{V} = \mathcal{F}$ contains the elementary
fault mechanisms and an edge $(A,B) \in \mathcal{E}$ indicates that
constituents $A$ and $B$ occur together in an observed compound
composition. The three context regimes can therefore be interpreted as
different graph-withholding operations: recombination removes only the
target edge $(A,B)$, one-sided context zero-shot removes all training edges
incident to one target constituent, and two-sided context zero-shot removes
all training edges incident to either target constituent.
Figure~\ref{fig:regimes} illustrates these operations on the
compound-composition subgraph.

\begin{figure*}[!t]
\centering
\includegraphics[width=\textwidth]{fig1_regimes}
\caption{Graph interpretation of the three compound-context evaluation
regimes. Nodes denote fault mechanisms participating in the compound
subgraph, edges denote observed compound compositions, and the red edge
denotes the target test composition. Recombination withholds only the
target edge; one-sided context zero-shot removes all compound-context
exposure for one target constituent; and two-sided context zero-shot
removes compound exposure for both target constituents.}
\label{fig:regimes}
\end{figure*}

Let the target test composition be
\begin{equation}
\mathcal{C}^{*} = \{A, B\}.
\end{equation}
Three levels of compound-context exposure were considered.

\subsubsection{Recombination}
The exact pair $A{+}B$ was excluded from training, while both $A$ and $B$
remained observable in other compound compositions. If
$d_{\mathrm{cmp}}(f)$ denotes the number of distinct compound contexts
containing constituent $f$ in the outer training set, the recombination
regime satisfies
\begin{equation}
d_{\mathrm{cmp}}(A) > 0, \qquad d_{\mathrm{cmp}}(B) > 0.
\end{equation}
Thus, the constituent mechanisms were familiar in compound operation, but
their particular edge or pairing was novel.

\subsubsection{One-Sided Compound-Context Zero-Shot}
In this regime, one constituent was completely removed from all compound
contexts while remaining available as a single fault. Without loss of
generality,
\begin{equation}
d_{\mathrm{cmp}}(A) = 0, \qquad d_{\mathrm{cmp}}(B) > 0.
\end{equation}
Both directional cases, $A$-unseen and $B$-unseen, were evaluated
separately.

\subsubsection{Two-Sided Compound-Context Zero-Shot}
The most restrictive context regime removed all compound examples
containing either constituent,
\begin{equation}
d_{\mathrm{cmp}}(A) = 0, \qquad d_{\mathrm{cmp}}(B) = 0,
\end{equation}
while preserving their eligible single-fault examples. Consequently,
neither constituent had been observed in any compound context before
testing on $A{+}B$.

Each composition-context regime was further crossed with operating-condition
exposure. Let $c_i$ denote the operating-condition identifier of run $i$.
For the seen-condition setting,
\begin{equation}
c_{\mathrm{test}} \in \mathcal{C}_{\mathrm{train}},
\end{equation}
whereas for the unseen-condition setting,
\begin{equation}
c_{\mathrm{test}} \notin \mathcal{C}_{\mathrm{train}}.
\end{equation}
The complete experimental design therefore consisted of
\begin{equation}
3 \ \text{compound-context regimes} \times 2 \ \text{condition regimes},
\end{equation}
resulting in six principal evaluation scenarios.

All splitting was performed at the experimental-run level before any
classifier fitting. Composition identities were defined from the
constituent vector rather than severity-specific class names, thereby
preventing severity variants of the same constituent combination from
unintentionally leaking across the composition boundary. Because either
member of a target compound could be deprived of prior compound-context
exposure, the one-sided regime was evaluated in both directions.
Consequently, this regime produced 216 prediction rows from 108 unique
physical compound runs, whereas the recombination and two-sided regimes
each produced 108 prediction rows from 108 unique runs. The two directional
predictions associated with the same physical run were never treated as
independent experimental units in uncertainty estimation; all bootstrap
analyses were clustered by physical run.

\subsection{Evaluation Metrics}
\label{sec:metrics}

Performance was evaluated using metrics that directly quantify recovery of
the constituent set rather than only conventional sample-level
classification. The overall constituent recall was defined as
\begin{equation}
R_{\mathrm{const}} =
\frac{\sum_{i=1}^{N}\sum_{k=1}^{K} y_{i,k}\hat{y}_{i,k}}
     {\sum_{i=1}^{N}\sum_{k=1}^{K} y_{i,k}}.
\end{equation}
Exact constituent-set recovery was computed as
\begin{equation}
\mathrm{EM} = \frac{1}{N}\sum_{i=1}^{N}
\mathbb{I}\left(\hat{\mathbf{y}}_i = \mathbf{y}_i\right).
\end{equation}
The average number of incorrectly added mechanisms per prediction row was
\begin{equation}
\mathrm{FA} = \frac{1}{N}\sum_{i=1}^{N}\sum_{k=1}^{K}
\left(1 - y_{i,k}\right)\hat{y}_{i,k}.
\end{equation}
The fraction of compound predictions for which at least one true
constituent was recovered was
\begin{equation}
R_{\mathrm{any}} = \frac{1}{N}\sum_{i=1}^{N}
\mathbb{I}\left(\sum_{k=1}^{K} y_{i,k}\hat{y}_{i,k} > 0\right),
\end{equation}
whereas recovery of all true constituents, irrespective of additional false
activations, was defined as
\begin{equation}
R_{\mathrm{all}} = \frac{1}{N}\sum_{i=1}^{N}
\mathbb{I}\left(\sum_{k=1}^{K} y_{i,k}\hat{y}_{i,k}
= \sum_{k=1}^{K} y_{i,k}\right).
\end{equation}
Micro-averaged $F_1$ was calculated over all constituent decisions as
\begin{equation}
\Fmicro = \frac{2TP}{2TP + FP + FN}.
\end{equation}
Because some evaluation subsets contained only a subset of the nine
constituents, macro-$F_1$ was additionally calculated over active
constituents,
\begin{equation}
\Factive = \frac{1}{|\mathcal{K}_{+}|}
\sum_{k \in \mathcal{K}_{+}} F_{1,k},
\end{equation}
where $\mathcal{K}_{+}$ denotes the constituents having at least one
positive test instance.

\subsection{Within-Pair Detectability and Cross-Partner Transfer}
\label{sec:transfer}

To determine whether poor constituent recovery reflected an absence of
measurable evidence, an oracle detectability analysis was performed. For
constituent $A$ in target compound $A{+}B$, the compound state was compared
directly with the corresponding single-$B$ state,
\begin{equation}
A{+}B \quad \text{vs.} \quad B.
\end{equation}
The resulting within-pair discriminability was quantified using ROC AUC and
denoted $\mathrm{AUC}^{\mathrm{within}}_{A|B}$.

A stronger transfer analysis evaluated whether evidence for $A$ learned
from eligible compound contexts other than the target composition remained
useful when $A$ occurred with partner $B$. The target pair $A{+}B$ was
excluded from training and the resulting fixed-orientation score was
evaluated on $A{+}B$ vs.\ $B$. The conventional directional cross-partner
AUC was denoted $\mathrm{AUC}^{A|B}_{\mathrm{dir}}$.

Because an AUC below 0.5 can reflect reversed ranking rather than absence
of separability, an orientation-free diagnostic measure was additionally
defined as
\begin{equation}
\mathrm{AUC}_{\mathrm{sep}} =
\max\left(\mathrm{AUC}_{\mathrm{dir}}, 1 - \mathrm{AUC}_{\mathrm{dir}}\right)
= 0.5 + \left|\mathrm{AUC}_{\mathrm{dir}} - 0.5\right|.
\label{eq:aucsep}
\end{equation}
$\mathrm{AUC}_{\mathrm{sep}}$ was used only to characterize rank
separability and was not interpreted as deployable classification
performance, because post-hoc inversion of an unknown target score would
require target-label information. The fixed-orientation transfer gap was
defined as
\begin{equation}
\Delta_{\mathrm{transfer}} =
\mathrm{AUC}^{\mathrm{within}}_{A|B} - \mathrm{AUC}^{A|B}_{\mathrm{dir}}.
\end{equation}
This quantity characterizes degradation of the learned score orientation
and should not be interpreted directly as loss of orientation-free
separability.

Because the released recordings were acquired across multiple recording
dates, an additional recording-sensitivity analysis was performed.
Acquisition timestamps were recovered from the original file names for 287
of the 288 runs, spanning seven recording dates. The only duplicated
fault--condition state available in two distinct recordings was used as a
nuisance-control experiment. In addition, compound pairs acquired on the
same date and sharing one constituent were compared under matched operating
conditions to reduce recording-date mismatch in the interpretation of
within-pair separability. These analyses were treated as sensitivity
controls rather than as causal estimates of acquisition-session effects.

\subsection{Physical-Signature Preservation Analysis}
\label{sec:preservation}

A complementary feature-level analysis characterized how predefined
physical signatures changed when a constituent was embedded in a compound
fault. Only the mechanism-specific curated feature families defined
independently of the cross-partner results were used. For feature $j$
associated with constituent $A$, let $\delta_j^{A}$ denote its signed
standardized evidence under the single-fault reference and
$\delta_j^{A|B}$ denote the evidence for the same feature when $A$ occurred
together with partner $B$. Feature preservation was summarized as
\begin{equation}
\rho_j^{A|B} = \frac{\delta_j^{A|B}}{\delta_j^{A}}.
\label{eq:preservation}
\end{equation}
Negative values indicate feature-level sign reversal, values between zero
and one indicate suppression, values around one indicate approximate
preservation, and values greater than one indicate amplification. This
analysis was used to characterize heterogeneity of predefined physical
signatures and was not interpreted as a causal decomposition of the
multivariate classifier score.

\subsection{Statistical Analysis}
\label{sec:stats}

Uncertainty of the principal constituent-level metrics was estimated using
nonparametric cluster bootstrap resampling at the physical-run level. In
the one-sided regime, both directional predictions associated with a
physical run were retained within the same bootstrap cluster. In addition
to protocol-specific confidence intervals, paired bootstrap contrasts were
computed because the same physical compound runs contributed to the
compared context regimes. Three prespecified contrasts were evaluated
separately under seen and unseen operating conditions:
\begin{align}
\Delta_{R-O} &= M_{\mathrm{recomb}} - M_{\mathrm{one}},\\
\Delta_{R-T} &= M_{\mathrm{recomb}} - M_{\mathrm{two}},\\
\Delta_{O-T} &= M_{\mathrm{one}} - M_{\mathrm{two}},
\end{align}
where $M$ denotes the metric of interest. The 95\% confidence interval was
obtained from the 2.5th and 97.5th percentiles of the paired bootstrap
distribution. A contrast was considered resolved when this interval
excluded zero. To formally test whether the effect of compound-context
deprivation changed under operating-condition shift, a
difference-of-differences statistic was also evaluated,
\begin{equation}
\Delta_{\mathrm{int}} = \Delta_{\mathrm{unseen}} - \Delta_{\mathrm{seen}},
\label{eq:interaction}
\end{equation}
where $\Delta$ denotes the paired difference between recombination and the
corresponding context-zero-shot regime. Confidence intervals were estimated
using 5000 paired cluster-bootstrap resamples at the physical-run level.
For the auxiliary closed-set experiment, 95\% confidence intervals were
estimated using 5000 stratified run-level bootstrap resamples within the
original 24 diagnostic classes.

For cross-partner transfer, evidence for reversed score orientation was
tested using
\begin{equation}
T = \left|\mathrm{AUC}_{\mathrm{dir}} - 0.5\right|.
\label{eq:polarity}
\end{equation}
Positive and negative scores were exchanged within matched operating
conditions under permutation, thereby preserving the paired experimental
structure. Family-wise multiplicity across the 18 directed
constituent--partner relations was controlled using the Holm procedure. A
relation was described as a statistically supported polarity reversal only
when its directional AUC was below 0.5, its bootstrap interval remained
below 0.5, and the Holm-adjusted permutation analysis supported departure
from chance orientation.
